\scanpage{300}
\deleted{on peut exiger \unsure{que} $f_1($}\inserted{d'où $g_{\unsure{1|2}}(t)=g_{\unsure{2|1}}(1/t)$} \unsure{et|en} \uline{$g_1$ \deleted{\unsure{a}}est une constante}. Donc les \uline{pts de $Y$ ne sont pas séparés}

On peut aussi passer au quotient par le groupe $\mathbb Z/2$ [\unsure{opérant|agissant} sur $X_1\amalg X_2$ par \unsure{symétrie|échange}, et sur $Y$ par $t\to1/t$] on trouve

% slanted
\margin{\unsure{annulé|annuler}, à \unsure{ré-pester|ré-percuter|rejeter}}

\diagram{\begin{tikzcd}[ampersand replacement=\&]
(X_1\amalg_Y X_2)/\mathbb Z_2\\
X_1\amalg X_2/\mathbb Z/2\simeq X_1\arrow[u]
\end{tikzcd}}

où $X_1$ est \unsure{q. affine}, \deleted{\unsure{comme}} mais non $X_1\amalg_Y X_2/\mathbb Z_2$ puisque $X_1\amalg_Y X_2$ le serait \unsure{aussi}.

Il faut cependant vérifier que tte \deleted{partie}\inserted{partie finie} \deleted{\unsure{trajectoire}} de $X_1\amalg_Y X_2$ par $\mathbb Z/2$ est contenue dans un ouvert affine : c'est facile. [prendre une \unsure{droite} $D$ de $\mathscr X$, \inserted{\unsure{distincte de l'axe}} \deleted{\unsure{par contre}} pour une partie finie, regarder $(\mathscr X-D)\amalg(\mathscr X-D)\unsure{\mathpunct{.}|=}$

% archive sheet number: 299

